\section{Formal Problem Statement and Objective Design}

This section formalizes the adaptive governance objective implemented in QAI-Chain. The purpose is to make explicit what is optimized, what is observed, and where approximation enters the stack.

\subsection{State, Action, and Transition Semantics}

At environment step $t$, the governance controller receives a state vector
\begin{equation}
\mathbf{s}_t = [d_t, \bar{\ell}_t, \bar{\tau}_t, \bar{m}_t, \bar{h}_t, \bar{u}_t] \in \mathbb{R}^6,
\end{equation}
where $d_t$ is chain difficulty, $\bar{\ell}_t$ is a latency proxy, $\bar{\tau}_t$ is throughput proxy, $\bar{m}_t$ is mempool pressure, $\bar{h}_t$ is hash-rate proxy, and $\bar{u}_t$ aggregates utilization-like telemetry. The action is scalar and bounded:
\begin{equation}
a_t \in [-1, 1], \qquad d_{t+1} = \operatorname{clip}(d_t + a_t, d_{\min}, d_{\max}).
\end{equation}
The transition is stochastic due to simulated workload and mining outcomes.

\subsection{Reward Decomposition}

The implemented reward follows a weighted decomposition:
\begin{equation}
r_t = w_{\tau}\,\tau_t + w_{\ell}\,\ell_t + w_{\delta}\,\delta_t - \lambda\,|d_t - d^\star|,
\end{equation}
where $d^\star$ is the target difficulty, $\tau_t$ is throughput utility, $\ell_t$ is latency utility (higher is better after sign normalization), and $\delta_t$ approximates decentralization utility. In our reference experiments we optimize episodic return
\begin{equation}
J(\pi) = \mathbb{E}_{\pi}\left[\sum_{t=0}^{T-1} \gamma^t r_t\right],
\end{equation}
with horizon $T=25$ and per-run episodic aggregation over 35 or 60 episodes depending on suite depth.

\subsection{Policy Optimization Objective}

The PPO branch uses a clipped surrogate objective \citep{schulman2017ppo}:
\begin{equation}
\mathcal{L}^{\text{clip}}(\theta) = \mathbb{E}_t\Big[\min\big(r_t(\theta)\hat{A}_t, \operatorname{clip}(r_t(\theta), 1-\epsilon, 1+\epsilon)\hat{A}_t\big)\Big],
\end{equation}
where
\begin{equation}
r_t(\theta) = \frac{\pi_\theta(a_t\mid s_t)}{\pi_{\theta_{\text{old}}}(a_t\mid s_t)}.
\end{equation}
The full loss is
\begin{equation}
\mathcal{L}(\theta, \phi) = -\mathcal{L}^{\text{clip}}(\theta) + c_v\,\mathcal{L}^{\text{value}}(\phi) - c_e\,\mathcal{H}[\pi_\theta],
\end{equation}
where $\mathcal{H}$ is policy entropy regularization and $c_e$ is entropy coefficient with linear annealing schedule.

\subsection{Advantage Normalization and Stability}

We compare PPO with and without advantage normalization:
\begin{equation}
\hat{A}_t \leftarrow \frac{\hat{A}_t - \mu_A}{\sigma_A + 10^{-8}}.
\end{equation}
Normalization can improve conditioning in mini-batch gradient steps by reducing scale variance across updates \citep{sutton2018reinforcement}. In our environment, removing it can increase average reward in some seeds, but with larger instability under policy drift.

\subsection{Complexity Notes}

For a policy/value MLP with width $h$, input dimension $d$, and scalar action head, forward complexity per step is $\mathcal{O}(dh + h^2)$. For the 6-qubit variational policy, quantum circuit evaluation dominates per-step inference cost. The runtime tradeoff in this work is accepted to gain uncertainty quality in high-switch governance regimes.

\subsection{Design Principle: Honest Baselines First}

We keep strong simple baselines. In constrained control surfaces, proportional heuristics can dominate RL policies when target structure is directly exposed. Reporting this outcome directly is necessary for scientific validity and empirical comparison \citep{pineau2021improving}.
